\documentclass[manuscript,screen,review]{jair}
\setcopyright{cc}
\acmDOI{}
\acmVolume{}
\JAIRAE{}
\JAIRTrack{}
\RequirePackage[datamodel=acmdatamodel,style=acmauthoryear,backend=biber,giveninits=true,uniquename=init]{biblatex}
\renewcommand*{\bibopenbracket}{(}
\renewcommand*{\bibclosebracket}{)}
\addbibresource{refs.bib}
\usepackage{booktabs}
\usepackage{graphicx}
\usepackage{amsmath}
\usepackage{url}
\usepackage{xcolor}
\usepackage{tabularx}
\usepackage{multirow}
\usepackage{placeins}
\input{a3_results/revision5_numbers.tex}
\begin{document}

\title[Same Trajectories, Different Rankings]{Same Trajectories, Different Rankings: Success Timing and Protocol Choices in Latent World-Model Planning}
\author{Ruiqi Wang}
\authornote{Corresponding author.}
\orcid{0009-0005-4942-5821}
\email{riw010@ucsd.edu}
\affiliation{%
  \institution{Department of Physics, University of California San Diego}
  \streetaddress{9500 Gilman Drive}
  \city{La Jolla}
  \state{CA}
  \postcode{92093}
  \country{USA}}
\renewcommand{\shortauthors}{Wang}

\begin{abstract}
\textbf{Background:}\quad  Recent latent-world-model papers report success on the four-task LeWM benchmark, usually counting an episode as successful when a trajectory first enters a target region.\par\noindent\textbf{Objectives:}\quad  We measure how success timing, observation history, planner settings, and sampling choices change the reported number and the method ranking.\par\noindent\textbf{Methods:}\quad  Public checkpoints and released code were evaluated on paired cases with three fixed planner random-seed streams. We changed one protocol choice at a time and applied any-step, block-boundary, final-step, and sustained-window labels to the same 50-step trajectories. The three preregistered predictions concern any-step success and were tested on 100 held-out cases per task; endpoint and sustained-window comparisons re-score the same trajectories.\par\noindent\textbf{Results:}\quad Three-frame history raises Reacher any-step success by \ReacherHistoryDeltaRevision{} percentage points and final-step success by \ReacherHistoryFinalCI{} points. Yet three-frame LeWM is successful at the final step in \ReacherLeWMThreeFinal\% of runs and at all ten final environment-step endpoints in only \RfourReacherThreeAllTen\%; endpoint occupancy is not sustained occupancy. AP and AD-WM retain common-protocol gains of \TwoRoomADGain{} to \CubeAPGain{} points on TwoRoom and Cube. A corrected continuation adapter passes checks on all \RfourFormalCount{} saved Reacher trajectories; our earlier adapter had allowed an unintended internal reset after success. Corrected comparisons show reversals for the same controller pair across arrival, endpoint, and sustained-window labels. In Reacher, three-frame LeWM minus AP is \RfourReacherThreeAPAny{} points for any-step, \RfourReacherThreeAPFinal{} for final-step, and \RfiveReacherThreeAPAllTen{} for all-last-ten success. In PushT, AP minus one-frame LeWM is \APLeWMPushTAnyCI{}, \APLeWMPushTFinalCI{}, and \RfourWindowSix{} points, respectively: the final-step and sustained-window intervals lie on opposite sides of zero.\par\noindent\textbf{Conclusions:}\quad First arrival, occupancy at designated endpoints, and sustained-window occupancy answer different control questions and can rank the same controllers differently. A benchmark result requires explicit check times, continuation semantics, history construction, planner budget, and paired sampling.

\end{abstract}
\keywords{world models, model predictive control, evaluation protocols, reproducibility, benchmark auditing}
\maketitle

\section{Introduction}
\IntroExampleSentence{} First-arrival success asks whether a controller reaches the target within the budget. Endpoint occupancy asks whether it is there at designated block boundaries or at the final step. Sustained-window occupancy asks how much, or whether all, of the final window is spent inside. These three questions produce different rankings from the same continuous trajectories.

Observation-and-action history has a large task-specific effect. In Reacher, three-frame LeWM gains \ReacherHistoryDeltaRevision{} points in any-step success and \ReacherHistoryFinalCI{} points at the final step relative to one-frame LeWM. Its final-step rate is \ReacherLeWMThreeFinal\%, but all-ten-step success is only \RfourReacherThreeAllTen\%. AP instead has \ReacherAPFinal\% final-step success, \RfourReacherAPOccTen\% mean last-ten occupancy, and \RfourReacherAPAllTen\% all-ten success. Thus endpoint improvement does not imply sustained holding (Table~\ref{tab:window-main}).

The same controller pair can reverse order across arrival, endpoint, and sustained-window labels. For three-frame LeWM minus AP in Reacher, any-step, final-step, and all-last-ten contrasts are \RfourReacherThreeAPAny{}, \RfourReacherThreeAPFinal{}, and \RfiveReacherThreeAPAllTen{} points; all three intervals exclude zero, with two opposite orderings. For AP minus one-frame LeWM in PushT, the corresponding contrasts are \APLeWMPushTAnyCI{}, \APLeWMPushTFinalCI{}, and \RfourWindowSix{} points. Any-step success is unresolved, while final-step and sustained-window success establish opposite orderings.

Common-protocol gains remain large for AP and AD-WM on TwoRoom and Cube. Interpreting post-success behavior also requires a continuous physical rollout. Continuing a Reacher rollout past success requires disabling the task's success termination. Our earlier continuation adapter cleared only the outer \texttt{terminated} flag; dm\_control had already scheduled a reset, so the arm was re-randomized on the next step. The corrected adapter preserves every saved any-step outcome in all \RfourFormalCount{} formal records.

This distinction matters because the LeWM benchmark is shared by several latent-planning papers, while the published descriptions expose different parts of the evaluation protocol. Our selected inventory contains \PublicSourceCount{} sources with results on at least one of the four tasks and \FourTaskSourceCount{} sources reporting all four (search cutoff: 2026-10-12 UTC). The records vary in how they specify success thresholds, early termination, observation history, CEM iterations, case counts, and planning seeds. These fields determine what a percentage measures and how finely a comparison can be resolved.

We ask three questions. Which choices move a measured success number? Which method differences survive when cases, streams, and planner settings are aligned? Does changing the success definition change the ranking? We answer them with a paired read-only evaluation of public checkpoints and code. Each case is run with three fixed planning streams, protocol factors are changed one at a time, and the same full-budget trajectories support first-arrival, designated-endpoint, and sustained-window measurements. The original any-step predictions use held-out confirmation.

The remaining question is empirical precision. Three planning streams over 100 cases have a typical paired half-width of \ReacherNoiseHalfwidth{} points in Reacher and \PushTNoiseHalfwidth{} points in PushT under the earlier any-step noise analysis. The continuous-run audit, method alignment, and uncertainty measurements together identify which evaluation choices change conclusions.

The paper contributes a paired evaluation design that exposes these protocol effects, a reproduction table with measured intervals, a full-trajectory comparison of arrival, endpoint, and sustained-window labels, a literature protocol inventory, and a reporting checklist. Section~2 introduces the planner and tasks; Section~3 describes the audit; Sections~4--8 present reproduction, factor effects, common-protocol gains, success-label reversals, and resolution; Sections~9--10 discuss implications and conclusions. The appendices provide the timeline, complete tables, planner-iteration analysis, compute accounting, and a reproducibility checklist.

\section{Background and Related Work}
\subsection{Latent world-model planning}
A latent world model maps an observation and an action into a learned representation and predicts future representations without reconstructing every pixel. A model-predictive-control (MPC) planner samples candidate action sequences, predicts their consequences, scores the candidates against a goal, and executes a prefix before planning again. The cross-entropy method (CEM) narrows the candidate distribution by repeatedly retaining an elite set and refitting its mean and variance. The LeWM evaluation executes a complete 25-step plan before the next query in its reference configuration; alternative released planners shorten the executed prefix or alter the goal score. Figure~\ref{fig:overview} depicts the objects that are held fixed and the labels applied after a full trajectory has been produced.

\begin{figure}[t]\centering\includegraphics[width=\linewidth]{figures/FIG1_AUDIT_OVERVIEW.png}\Description{Paired evaluation design and a continuous Reacher example with joint-error curves and success labels.}\caption{Measurement times and an example trajectory. A case-stream identifies a paired trajectory; the case is the bootstrap unit. The right panel shows a corrected continuous Reacher trajectory, its two absolute joint errors and the 0.05 rad threshold. The left panel marks block checks, 25/50-step plan endpoints, AP re-decisions, window reads, and the Reacher OR/END microstep distinction. The right panel is the recorded corrected trajectory.}\label{fig:overview}\end{figure}

This paper evaluates the original LeWM checkpoint, Fast-LeWM, AP (Anchored Planning), and AD-WM. Fast-LeWM uses a one-frame input. AP anchors short plans to a frozen world model and retrieval bank. AD-WM changes the action-conditioned objective and reports three training seeds. These methods share the task interface but expose different planner choices, which makes protocol alignment necessary before comparing their percentages \parencite{lewm2026,fastlewm2026,ap2026,adwm2026}.

\subsection{Tasks and success tests}
The four tasks cover a two-joint arm, a planar pusher, a point agent crossing two rooms, and a three-dimensional cube manipulation environment. The audit uses the task thresholds in the released interfaces; no threshold is changed. Table~\ref{tab:tasks} summarizes the tests and the common 50-step budget.
\input{tables/table_tasks.tex}

\subsection{Published protocols and released configurations}
LeWM Appendix D describes 30 CEM iterations for PushT and 10 for the other tasks; Appendix B states 30 without a task qualification. The released evaluation configuration uses 30 throughout. The appendix lists three history frames for PushT and Cube, one for TwoRoom, and does not state a Reacher value; the library default and the released configuration supply one frame. Case counts and planning seeds also differ across sources. Table~\ref{tab:literature} records the value or “not specified” for each source rather than filling gaps by inference.
\input{tables/table_protocol_literature.tex}

The selected inventory contains \PublicSourceCount{} sources, including benchmark papers, planner papers, and objective studies. \FourTaskSourceCount{} cover all four tasks. We included a source when it reported a result, protocol, or released configuration for at least one benchmark task and was publicly identifiable by the 2026-10-12 UTC cutoff; the table is an inventory of selected sources. Evidence labels distinguish a value written in the paper (paper), determined from a released configuration (config), determined from source code (code), or not determinable (not stated). The protocol fields are not interchangeable: an unspecified history value cannot be treated as a reported three-frame setting, and a paper's case count does not define the number of streams in a released configuration. We therefore reproduce source-specific numbers in Section~4 and use the reference protocol for aligned contrasts.

\subsection{Related work}
Evaluation methodology in reinforcement learning has documented the influence of random seeds, sample size, environment versions, and reporting summaries on conclusions \parencite{henderson2018,machado2018,agarwal2021,gundersen2024}. Recent work on benchmark reporting and multi-task suites extends this concern to task definitions and environment wrappers \parencite{patterson2024,metaworldplus2025}. Our audit instantiates those concerns for latent world-model planning by holding cases and planning streams paired within each comparison.

Several recent papers study the same benchmark. PPWM is the closest predecessor: it evaluates the same four tasks with LeWM, Fast-LeWM, PP-LeWM, and several planners, using 100 episodes, three planner seeds, and complete 50-step traces to report both anytime and terminal success; its Reacher table already shows high anytime and lower terminal rates and therefore contains the ingredients for metric-dependent rankings \parencite{ppwm2026}. Singh's independent reproduction is a direct protocol-audit predecessor: it reproduces the TwoRoom task and checks how the evaluation convention determines the result \parencite{singh2026}; a companion paper studies the objective bottleneck \parencite{singhobjective2026}. AP's Appendix D.2 separately audits the LeWM configuration and evaluation protocol \parencite{ap2026}. Vesessook and Yang study planner scoring times and replanning intervals \parencite{vesessook2026}; ProWorld discusses progress-aware scoring in long-horizon goal reaching \parencite{proworld2026}. AD-WM reports Cube results under its native setup and compares Fast-LeWM, Sub-JEPA, and INTACT public checkpoints \parencite{adwm2026}. We extend these lines with a protocol audit across later methods, fixed case-and-stream paired effects, preregistered held-out confirmation, and an explicit check of continuation semantics after success.

A second line of work changes the planning objective rather than the evaluation. TEMPO, TRM, and decision-metric alignment methods improve the match between a latent objective and the action that is eventually executed \parencite{tempo2026,trm2026,decisionmetric2026}. In other settings, success definitions also change interpretation: dm\_control's original Reacher task gives reward at each step spent inside the target, so episode return measures target occupancy more closely than first entry \parencite{dmcontrol2018}. VIGIL contrasts a world state that is complete at episode end with an agent that correctly reports completion \parencite{vigil2026}. These examples motivate recording check times and termination rules alongside a success percentage.

\section{Audit Design}
\subsection{Checkpoints, code, and cases}
We evaluated the public LeWM checkpoint, Fast-LeWM, AP, and AD-WM with the released adapters and task interfaces. The adapter and checkpoint identifiers were fixed before the formal runs. Fast-LeWM is evaluated with one-frame history because its released interface does not provide a three-frame condition. AD-WM's three training seeds are averaged within case and planning stream before contrasts are formed. Methods without public weights and code were inventoried but not substituted with private implementations.

The selected public methods were fixed by the earlier decision record: a public checkpoint and adapter available by the freeze date, a reproducible interface for at least one task, a distinct planner or objective choice, and a budget compatible with the paired audit. This is why LeWM, Fast-LeWM, Anchored Planning, and AD-WM are evaluated here. EpiWM and other public-weight candidates were considered in the earlier candidate inventory but are not part of the evaluated set when their released interface, checkpoint identity, or task coverage did not meet all four criteria.

The task interfaces come from the stable-worldmodel platform \parencite{stableworldmodel2026}, audited at commit \texttt{63988116}. Success termination is a legitimate task design; the error was in our earlier post-success continuation adapter. The released first-arrival evaluation never continues past success and is unaffected. For Reacher, the corrected adapter preserves the released strict per-joint condition $|q_j-q_j^*|<0.05$ rad but prevents success from terminating the underlying dm\_control episode during the 50-step budget. Both native microsteps are executed for every environment step. We record the original predicate at each microstep, the two step labels, physical time, joint position and velocity, target configuration, and actual reset events. Appendix~\ref{app:reacher-verification} reports the runtime checks and final-code technical gates.

The main case set contains 200 cases per task and the held-out case set contains 100 cases per task. Each case comes from a different episode and contributes one evaluation window; no episode contributes two cases to the same cohort. The original-main and fresh-main cohorts are each 100 cases per task. The held-out set excludes whole episodes used by this study's earlier evaluations, retrieval bank, and episodes used in the author's earlier retraining experiments on the same benchmark (those training results are not reported here). The case list is sampled from the benchmark pool with the recorded target-distance/start-state selection procedure, rather than selected after inspecting success. This exclusion prevents analysis choices in the study from selecting the held-out episodes.
\input{tables/table_cohorts_revision.tex}

\subsection{Native settings and study interfaces}
The public methods expose different planner interfaces. Table~\ref{tab:method-mapping} maps each native setting to the common study interface. Native reset restores the task's released state variables (Reacher joint position/velocity, PushT state, TwoRoom proprioception, or Cube joint and block state); native goal uses the task's released goal state, with the Cube target reconstructed from the recorded target position and quaternion. Goal images are dataset endpoint images when the interface supplies them. Actions are converted through the released action normalisation and clipped at the environment bounds. The common study executes the same 50 raw environment steps and continues after an any-step success label; the corrected Reacher implementation prevents success from triggering an internal reset.
\input{tables/table_method_mapping_revision.tex}

The planner and environment clocks use raw environment steps. A model prediction may cover a longer imagined horizon, but the released LeWM interface emits and executes a 25-step plan comprising five five-step action blocks, then queries again if budget remains. Steps 5, 10, \ldots, 50 are action-block boundaries; 25 and 50 are endpoints of LeWM plans. AP re-decides every five raw steps and re-retrieves at that boundary. An action-block boundary is therefore distinct from a re-decision boundary: LeWM re-decides after 25 steps, whereas AP re-decides after five. Three-frame history means the current observation together with the real observations at $t-5$ and $t-10$ and the two intervening executed action blocks; the first query has only the current frame. No zero padding or predicted frames are inserted.

\subsection{Paired streams and aggregation}
A case is the statistical unit. Every method in a contrast receives the same case, the same reset, and one of three fixed planner random-seed streams. Pairing removes case difficulty from the first-order comparison and makes the interval describe the remaining case-level variation. For AD-WM, training seeds are averaged within case before the planning-stream average. All formal traces run the full 50-step budget even when an any-step label would already be true.

\subsection{Protocol factors}
The reference protocol uses native reset, one-frame history, 30 CEM iterations, 300 candidates, 30 elites, five five-step blocks, and a full 50-step trajectory. We change one interface at a time where a released implementation permits it. Figure~\ref{fig:factors} shows history length, the CEM elite standard-deviation correction, goal offset, goal image, and native replanning effects on the mixed 200-case set. The separate ten-versus-thirty iteration comparison uses a held-out set and appears in Appendix~C.
\begin{figure}[t]\centering\includegraphics[width=\linewidth]{figures/FIG2_FACTOR_EFFECTS.png}\Description{Paired effects of protocol factors in Reacher, PushT, TwoRoom, and Cube.}\caption{Protocol-factor effects relative to the reference protocol. Error bars are case-bootstrap 95\% intervals; the figure uses the same paired cases and streams for each factor. Gray bands are the minimum--maximum numeric gains in the selected literature inventory, included only as a magnitude reference; their bounds are the inventory extrema, not a pooled estimate.}\label{fig:factors}\end{figure}

\subsection{Success labels and occupancy}
Let $e_t$ denote the task predicate at the end of environment step $t$, and let $a_t$ denote whether it holds at any success check within that step. Reacher executes two microsteps, so $a_t$ is their logical OR and $e_t$ is the second microstep predicate; the other tasks use their released step-end check. Any-step success is $\mathbb{1}(\max_{1\leq t\leq50}a_t=1)$, block-boundary success is $\mathbb{1}(\max_{t\in\{5,10,\ldots,50\}}e_t=1)$, and final-step success is $e_{50}$. Occupancy is $100\times\Pr(e_t=1)$, with training-checkpoint and planning-stream averaging inside each case. For the last $K$ steps, mean occupancy is $K^{-1}\sum_{t=51-K}^{50}e_t$ and all-window success is $\prod_{t=51-K}^{50}e_t$, with $K\in\{5,10\}$. They are computed per physical trajectory before averaging. Arrival, endpoint, and sustained-window labels address three control questions; block-boundary and final-step labels are two choices of endpoint checks. The OR/end distinction is quantified in Appendix~\ref{app:reacher-verification}.

\subsection{Statistics and registration}
For method $m$, checkpoint $k$, stream $s$, case $i$, and success definition $d$, let $Z_{ikmsd}\in[0,1]$ denote the binary label or mean window occupancy computed on one complete physical trajectory. Define $Y_{imsd}=K_m^{-1}\sum_{k=1}^{K_m}Z_{ikmsd}\in[0,1]$, where $K_m=3$ for AD-WM and $K_m=1$ otherwise. A nonlinear trajectory metric is calculated before checkpoint averaging; the averaged labels are not thresholded again. The estimand is the fixed-checkpoint, fixed-three-stream average over the case distribution, with the planning-stream mean formed inside each case:
\[
\widehat{\Delta}_{m,d}=\frac{1}{n}\sum_{i=1}^{n}\left[\frac{1}{S}\sum_{s=1}^{S}Y_{imsd}-\frac{1}{S}\sum_{s=1}^{S}Y_{i\mathrm{ref}sd}\right].
\]
The 95\% percentile interval resamples the same case indices for every method, stream, checkpoint, and success definition in each bootstrap replicate. AD-WM checkpoint seeds are averaged within case and stream before the contrast. This estimates the average effect for the fixed public checkpoints, fixed planner streams, and fixed protocol; it does not estimate retraining, a different stream set, a different retrieval bank, or another environment version. Contrasts use 5,000 case-bootstrap replicates. Degenerate cells also receive exact binomial intervals, while multi-stream or multi-checkpoint case means are never treated as independent binomial trials. The audit records source-specific reproduction, factor contrasts, arrival/endpoint/window labels, occupancy, rank correlations, and held-out prediction tests. The protocol sequence began with source reproduction, then native reset and held-out confirmation; after the scope was fixed, the held-out work retained the history, common-protocol, and success-definition questions and included the CEM-iteration comparison. The complete chronology and scope decisions are in Appendix~A.

The original preregistered any-step predictions retain confirmatory status under their original cases, configurations, and rules: the full saved-record audit finds zero old/new any-step outcome mismatches (Appendix~\ref{app:reacher-verification}). The correction registration froze the fix and analysis plan before the valid formal rerun (Table~\ref{tab:reacher-correction-timeline}). CEM-iteration comparisons, corrected endpoint effects, sustained-window results, and ranking changes are descriptive analyses, not a new held-out confirmation: those cases had already been seen.

\section{Can the Published Numbers Be Reproduced?}
Table~\ref{tab:replication} gives the source value, our measured value, its interval, and whether the published value lies within the preregistered bootstrap interval. Ten of 12 cells match. The Fast-LeWM PushT value is \ReproFastPushTOur\% versus a published \ReproFastPushTPaper\%. AP Reacher is \ReproAPReacherOur\% versus \ReproAPReacherPaper\% in the source record. Both use the released cases and query list. In the AP list, 30 source failures become local successes and two source successes become failures; the record does not identify a single cause.
\input{tables/table_replication_revision.tex}

\section{Which Unreported Choices Move the Number?}
\subsection{Factor effects relative to published gains}
Figure~\ref{fig:factors} places each paired factor effect on the four tasks; the corresponding intervals are in Table~\ref{tab:factors}. The history change is the largest stable factor in Reacher, while native replanning lowers Reacher by \ReplanReacherDelta{} points and PushT by \ReplanPushTDelta{} points. The re-planning comparison includes the native every-five-step schedule and warm start; the experiment did not separate those implementation components.
\input{tables/table_factor_effects.tex}

\subsection{History and software behavior}
Three-frame history adds \ReacherHistoryDeltaRevision{} points in Reacher any-step success (95\% CI \ReacherHistoryCIRevision), with changes of \PushTHistoryDeltaRevision{}, \TwoRoomHistoryDeltaRevision{}, and \CubeHistoryDeltaRevision{} points in the other tasks. Corrected final-step success rises from \ReacherLeWMOneFinal\% to \ReacherLeWMThreeFinal\%, a paired \ReacherHistoryFinalCI{}-point effect. The intervention provides observation and action history, and this study measures its effects without attributing them to a particular mechanism. One-frame any-step success is already \ReacherLeWMOneAny\%, leaving only \APReacherAnyDelta{} points below its ceiling.

The first query contains only the current frame. Across all \RfourDPairs{} paired case-streams, the first 25 raw actions are bitwise identical for one- and three-frame LeWM, with maximum absolute difference zero; their first action divergence is at step 26. First-plan OR and END labels also agree. Thus the first plan is identical by construction and verification, and both gains arise in the second plan, at steps 26--50. Table~\ref{tab:window-main} shows why these arrival and endpoint improvements should not be called sustained holding: three-frame LeWM's all-ten success is only \RfourReacherThreeAllTen\%.

The library comparison changes the CEM elite standard-deviation correction. Its paired effects are about \LibraryReacherDelta{} points in Reacher and \LibraryCubeDelta{} points in Cube, while PushT and TwoRoom intervals include zero. The library change and elite-correction intervention are numerically identical and are merged into one row because this correction is the substantive version difference. CEM 10 versus 30 iterations changes the measured rate by at most about \CEMMaxDelta{} points in these cells; the full iteration table and plot are in Appendix~C. Goal offset and goal image changes are smaller than the history and replanning effects.

The Cube goal-image factor has no success-rate estimate: one isolated goal-render equality check failed before any model call or environment step, so that condition was excluded. The failure does not remove or alter the other Cube factor conditions, whose input images and trajectories passed their recorded checks.

\section{How Much of the Published Gain Survives a Common Protocol?}
Figure~\ref{fig:common} compares each method with LeWM one-frame on the same held-out cases and streams. AP and AD-WM gains on TwoRoom are \TwoRoomAPGain{} and \TwoRoomADGain{} points; on Cube they are \CubeAPGain{} and \CubeADGain{} points. Their intervals exclude zero. In Reacher, three-frame history itself is \ReacherHistoryDeltaRevision{} points above one-frame LeWM; it is \ReacherThreeVsFast{} points above Fast-LeWM and \ReacherThreeVsAD{} points above AD-WM. \ReacherAPHistoryComparisonSentence{} PushT contrasts are within a few points.
\begin{figure}[t]\centering\includegraphics[width=\linewidth]{figures/FIG3_COMMON_PROTOCOL_GAINS.png}\Description{Paired method differences relative to one-frame LeWM on the held-out case set.}\caption{Common-protocol paired gains under any-step success. Points are contrasts and bars are case-bootstrap 95\% intervals.}\label{fig:common}\end{figure}
\input{tables/table_predictions_revision.tex}

\CorrectedPredictionSummary{} (Table~\ref{tab:predictions}). The main, fresh, and held-out cohorts show the same direction for the history and TwoRoom/Cube contrasts. The AP retrieval bank differs between the two evaluation rounds because all earlier evaluation episodes were removed before held-out selection; the cohorts are therefore reported separately.

\section{What Counts as Success?}
\subsection{Arrival and endpoint labels on one trajectory}
Table~\ref{tab:success} applies any-step, block-boundary, and final-step labels to identical 50-step traces. LeWM Reacher falls from \ReacherLeWMOneAny\% to \ReacherLeWMOneFinal\% when the label changes from any-step to final-step. AP falls from \ReacherAPAny\% to \ReacherAPFinal\%. The paired AP--LeWM difference is \APLeWMReacherAnyCI{} points under any-step and \APLeWMReacherFinalCI{} points under final-step; the same bootstrap gives a \APLeWMReacherChangeCI{}-point change in the AP--LeWM difference when the success definition changes. Figure~\ref{fig:success} shows all four tasks and all five configurations.
\input{tables/table_success_rates_revision.tex}
\begin{figure}[t]\centering\includegraphics[width=\linewidth]{figures/FIG4_SUCCESS_RATES_REV6.png}\Description{Any-step, block-boundary, and final-step rates for five configurations on four tasks.}\caption{Arrival (any-step) and endpoint (block-boundary and final-step) labels applied to the same held-out trajectories. Bar textures distinguish the three check-time choices in grayscale.}\label{fig:success}\end{figure}

\subsection{Correcting our continuation adapter}
Continuing a Reacher rollout past success requires disabling the task's success termination. Our earlier continuation adapter cleared only the outer \texttt{terminated} flag; dm\_control had already scheduled a reset, so the arm was re-randomized on the next step. The released first-arrival evaluation never continues past success and is unaffected. \ReacherRuntimeReplaySentence{} The corrected adapter retains the same predicate and dynamics while disabling this success termination.

The old/new comparison combines reset removal and the explicit END convention; the OR/END table in Appendix~\ref{app:reacher-verification} isolates label choice on the corrected trajectories. Any-step rates are identical before and after correction in every configuration, as intended, and the full \RfourFormalCount{}-record audit confirms zero paired any-step outcome mismatches.
\input{tables/table_reacher_termination_correction.tex}
Appendix~\ref{app:reacher-verification} (Table~\ref{tab:reacher-correction-contrasts}) compares the AP--LeWM contrasts and rank correlations under the two adapters.


Table~\ref{tab:code-audit-r7} traces the conditions behind this failure mode in public code. Success termination, an automatic reset, and continued sampling must coincide for post-success observations to cross an episode boundary; clearing an outer vector flag does not clear dm\_control\textquotesingle s internal reset state. Stopping at success or explicitly retaining episode boundaries avoids that interpretation; the table reports static paths, not a re-evaluation of published scores.

\subsection{Ranking reversals}
The same controller pair can reverse order across arrival, endpoint, and sustained-window labels. For three-frame LeWM minus AP in Reacher, any-step, final-step, and all-last-ten contrasts are \RfourReacherThreeAPAny{}, \RfourReacherThreeAPFinal{}, and \RfiveReacherThreeAPAllTen{} points; all three intervals exclude zero, with two opposite orderings. For AP minus one-frame LeWM in PushT, the corresponding contrasts are \APLeWMPushTAnyCI{}, \APLeWMPushTFinalCI{}, and \RfourWindowSix{} points. Any-step success is unresolved, while final-step and sustained-window success establish opposite orderings. Figure~\ref{fig:reversal} uses one-frame LeWM as the reference and shows arrival, endpoint, and sustained-window labels for every configuration. For PushT, the final-minus-any contrast change is \APLeWMPushTChangeCI{} points. TwoRoom\textquotesingle s any-step and final-step rates and ordering are close.

\begin{figure}[t]\centering\includegraphics[width=0.95\linewidth]{figures/FIG5_RANKING_REVERSALS.png}\Description{Method differences from one-frame LeWM under any-step, final-step, and all-last-ten success.}\caption{Arrival (any-step), endpoint (final-step), and sustained-window (all-last-ten) labels can order the same controllers differently.}\label{fig:reversal}\end{figure}
\input{tables/table_rankings_revision.tex}

Kendall correlations summarize the same result. Reacher's any-step versus final-step value is \TauReacherAnyFinal; PushT is \TauPushtAnyFinal; TwoRoom is \TauTworoomAnyFinal; and Cube is \TauCubeAnyFinal. The block-boundary correlations are in Table~\ref{tab:rankings}. A benchmark therefore has a metric choice in addition to a checkpoint choice.

\paragraph{Sustained-window rankings.}
Table~\ref{tab:window-main} adds the third question to first arrival and endpoint occupancy. In Reacher, three-frame LeWM leads AP at the final step, whereas AP leads in mean last-ten occupancy by \RfourWindowOne{} points and in all-ten success by \RfourWindowTwo{}. AP also exceeds one-frame LeWM by \RfourWindowThree{} and \RfourWindowFour{} points, respectively. PushT AP and Fast-LeWM both exceed one-frame LeWM on the corresponding window measures; all paired intervals are in Appendix~D. Sustained-window occupancy therefore supplies another ranking rather than a synonym for final-step success.
\input{tables/table_window_main.tex}

Table~\ref{tab:sensitivity-r7} checks threshold and window sensitivity: the Reacher final-versus-all-window reversal retains nonzero intervals at 0.05 and 0.08 rad for $K=5,10,15,20$; at 0.03 rad the all-window interval first touches zero at $K=10$, with zero contrasts at $K=15,20$. PushT AP retains the reversal for all four windows at the recorded threshold, whereas Fast-LeWM leads under both labels; the AP mean-occupancy interval includes zero at $K=5$.

\subsection{Occupancy over time}
Figure~\ref{fig:occupancy} reports the mean step-end predicate at each raw step, with checkpoint and stream averaging inside each case. \ReacherOccupancySentence{} PushT, especially LeWM and AP, peaks near step 25, declines during steps 26--49, and rises again toward step 50; Fast-LeWM has peaks at both endpoints. TwoRoom mostly peaks at step 50, while Cube peaks at steps 25--26 and declines before a smaller endpoint rebound. The planned endpoints therefore coincide with many arrivals, while mid-plan exits create the gap between any-step and final-step labels. The exact peak and middle-minimum values are listed in Table~\ref{tab:occupancy-peaks}. AP's Reacher curve has a five-step sawtooth, with peaks at re-decision boundaries; LeWM and AD-WM peak at steps 25 and 50. These curves describe success timing on the environment side; \textcite{vesessook2026} study the planner's scoring side.
\begin{figure}[t]\centering\includegraphics[width=\linewidth]{figures/FIG6_OCCUPANCY_REV6.png}\Description{Target occupancy over raw steps for Reacher, PushT, TwoRoom, and Cube.}\caption{Occupancy is the mean per-step predicate, averaging checkpoint and planning-stream indicators within each case; vertical guides mark steps 25 and 50. Each task has its own vertical scale; Reacher uses the corrected step-end predicate.}\label{fig:occupancy}\end{figure}
\input{tables/table_occupancy_peaks.tex}

\section{How Small a Difference Can One Evaluation Resolve?}
A 50-case, one-stream Reacher evaluation has a typical paired interval half-width near \NoiseReacherFifty{} points. With 100 cases and one stream it is \NoiseReacherHundred{} points; with three streams it is \NoiseReacherHundredThree{} points. PushT gives \NoisePushTFifty{}, \NoisePushTHundred{}, and \NoisePushTHundredThree{} points for the same three designs. These half-widths use any-step labels in the earlier Step-2 200-case cohort and are unaffected by the continuation correction. These ``typical'' values are medians of the paired bootstrap half-widths for the recorded 50-case subsets, with one or three streams averaged within each case; they are empirical resolution summaries, not minimum detectable effects. The variance of a paired difference is driven by the fraction of cases on which two methods disagree, so it is not determined by case count alone. The 100-case, three-stream design resolves broad gains such as those on Cube but cannot distinguish many one- or two-point claims. Figure~\ref{fig:noise} shows the comparison.
\begin{figure}[t]\centering\includegraphics[width=0.8\linewidth]{figures/FIG7_NOISE_SCALE.png}\Description{Typical paired interval half-width under three case and stream designs.}\caption{Resolution improves with both cases and planning streams.}\label{fig:noise}\end{figure}
The broader Step-0 candidate inventory (supplied in the materials) contains \LiteratureGainN{} numeric method--task gain claims on these four tasks. \LiteratureGainWithinN{} (\LiteratureGainWithinPct\%) have absolute magnitude at most 10 percentage points, the approximate resolution of a 50-case one-stream Reacher comparison. Examples include Fast-LeWM gains of +2 points on Reacher and 0 on PushT, and the LEAP PushT gain of +2 points. These are inventory values, reported without reinterpreting their protocols or assessing significance in any source paper.

\section{Discussion}
\subsection{What each ruler measures}
Any-step success asks whether the controller can bring the system through the target before the budget ends. Final-step success asks whether the system is in the target at a specified later time. Block-boundary success asks the same question at the five-step action-block boundaries. Sustained-window occupancy asks how much or whether all of a final window is spent in the target; it distinguishes endpoint arrival from remaining there throughout the window. These are different control questions. Methods in this benchmark are commonly tuned and reported with the any-step convention, so a paper should report final-step success alongside the conventional number. The step-wise reward in the original dm\_control Reacher task illustrates the same distinction: transient proximity can contribute reward even when the terminal configuration differs.

\subsection{A reporting checklist}
The minimum fields and observed scales are in Table~\ref{tab:reporting}. A paper should state the task threshold, all check times, early termination, history construction, CEM candidates/elites/iterations, horizon, replanning interval, software commit and dependency versions, case selection, planning streams, training seeds, aggregation order, and uncertainty interval. An occupancy curve or first-hit diagnostic allows readers to interpret an any-step value alongside the final state.
\input{tables/table_reporting.tex}

\subsection{Limitations}
The empirical comparisons cover one world-model family, four simulated tasks, three public method families, and a finite set of planning streams. The static continuation audit additionally traces the public code paths in Table~\ref{tab:code-audit-r7}; that source-code coverage is distinct from empirical validation on those systems. The held-out exclusion concerns this study's own case and retrieval choices; public pretraining exposure is unknown. The AP retrieval bank differs between the two evaluation rounds. Main-result thresholds are inherited from the released environments; Appendix~D varies the Reacher threshold, while missing PushT physical states prevent its threshold re-scoring. The causes of leaving a target and of the native replanning decrease are not separated. AP Reacher reproduction was \ReproAPReacherOur\% versus the source's \ReproAPReacherPaper\%; this unexplained first-arrival discrepancy is unrelated to the continuation reset. AP conclusions here concern the implemented and recorded configuration. Results may differ for other world-model families, planners, or physical robots.

\section{Conclusion}
Evaluation choices determine what a latent-planning benchmark score means. In these four tasks, history length materially changes Reacher's any-step success, common-protocol gains on TwoRoom and Cube remain large, and first-arrival, endpoint and sustained-window labels can change the ordering of the same controllers. The Reacher audit also shows why a post-success continuation adapter must disable the task's success termination. Reporting the success timing, history, planner budget, case selection, and planning streams with every percentage makes such comparisons auditable and scientifically interpretable.

\section*{CRediT authorship contribution statement}
Ruiqi Wang: Conceptualization, Methodology, Software, Validation, Formal analysis, Investigation, Data curation, Writing -- initial manuscript, Writing -- review and editing, Visualization, Project administration.

\section*{Declaration of competing interest}
The author declares that there are no known competing financial interests or personal relationships that could have appeared to influence the work reported in this paper.

\section*{Funding}
This research did not receive any specific grant from funding agencies in the public, commercial, or not-for-profit sectors.

\section*{Data availability}
Audit records and code are archived at Zenodo, DOI \href{https://doi.org/10.5281/zenodo.23350373}{10.5281/zenodo.23350373}, with records under CC BY 4.0 and audit code under MIT. The code mirror is \url{https://github.com/2GyounnnG/lewm-protocol-audit}. Public checkpoints and input datasets are identified by source URLs and SHA-256 hashes and are not redistributed.

\section*{Generative AI statement}
During the preparation of this manuscript, the author used Claude (Anthropic) and OpenAI Codex (OpenAI) to assist with the design of analyses, writing and debugging experiment, analysis, and plotting code, preparing and editing text, and checking the manuscript against journal guidelines. The author reviewed and edited the manuscript and takes full responsibility for its content.

\appendix
\FloatBarrier
\section{Protocol timeline and scope}
The evaluation was organized as a sequence of source reproduction, protocol-factor checks, stream-matched post hoc analyses, and held-out confirmation. PushT exact-start override was withdrawn and 21 affected conditions were rerun with native reset; 572 earlier partial records were excluded. The held-out definition excludes the study's prior cases and retrieval bank, plus episodes from the author's earlier retraining experiments on this benchmark, whose results are not reported here. Before held-out results existed, the replanning-mechanism grid and AP/AD-WM three-frame conditions were removed and a ten-versus-thirty CEM iteration comparison was added. Adapter-fidelity checks were carried over only for unchanged adapters; the common evaluation environment was rebuilt. Table~\ref{tab:timeline-revision} gives UTC times and document hashes.
\input{tables/table_timeline_revision.tex}

\FloatBarrier
\section{Complete factor and method results}
Table~\ref{tab:success} gives the full held-out rate table used in the main figures. The accompanying machine-readable CSVs retain successful-trajectory counts, case-majority rates, and interval endpoints for every task, configuration, and definition. Factor-level rows appear in Table~\ref{tab:factors}; source protocol rows appear in Table~\ref{tab:literature}.
\input{tables/table_ap_lewm_intervals.tex}
The AP--LeWM intervals in Table~\ref{tab:ap-lewm-intervals} use one case-bootstrap index vector for both labels and for the reported change. Reacher rows use corrected continuous trajectories and the sealed bootstrap seed; the original paired analyses for the other tasks are retained.
\input{tables/table_adwm_stream_revision.tex}
The AD-WM rows retain each public checkpoint and each of the three planner streams rather than treating them as independent cases. The machine-readable audit contains the same fields for the other evaluated methods; Reacher rows use the corrected continuation semantics.

\FloatBarrier
\section{CEM iteration comparison}
Figure~\ref{fig:cem} reports its paired differences for any-step, block-boundary, and final-step labels; CEM settings, history, cases, and streams are held fixed within each point.
\begin{figure}[htbp]\centering\includegraphics[width=0.9\linewidth]{figures/FIG_APPENDIX_CEM_ITERATIONS.png}\Description{Paired differences between ten and thirty CEM iterations.}\caption{CEM iteration comparison.}\label{fig:cem}\end{figure}
\input{tables/table_cem_numeric.tex}

\FloatBarrier
\section{Supplementary terminal-window metrics}
The final-$K$ occupancy, all-$K$ success, and post-first-hit exit metrics are summaries of the same continuous trajectories. These are additional trajectory-level summaries computed with the unchanged task predicate; they were not part of the preregistered predictions. Metrics are calculated on each checkpoint trajectory before averaging streams and training seeds. A first hit at the last step has no subsequent observation and is excluded from the conditional exit denominator. The tables for all tasks use this same aggregation rule.
\input{tables/table_terminal_window.tex}
\input{tables/table_window_intervals.tex}
\input{tables/table_sensitivity_revision7.tex}

\FloatBarrier
\section{JAIR reproducibility checklist}
\input{tables/table_jair_checklist_official.tex}

\FloatBarrier
\section{Reacher runtime verification and continuation code audit}\label{app:reacher-verification}
\input{tables/table_code_audit_revision7.tex}
\input{tables/table_formal_coverage.tex}
\input{tables/table_reacher_runtime_checks.tex}
The TECH replay covers 108 trajectories, not the full formal grid. It instantiates the uninstrumented native stable-worldmodel Reacher wrapper, bypasses the recording and continuation wrapper, disables only success termination, and sends each saved action through two native dm\_control steps. Comparisons include \texttt{qpos}, \texttt{qvel}, physical time and \texttt{target\_qpos}, before and after every microstep. The maximum discrepancy is zero with absolute tolerance $10^{-12}$. Predicates are separately recomputed from saved joint and target states. This is an independent replay path through the same simulator and dynamics, not an independent physics implementation.

All \RfourFormalCount{} formal records additionally undergo a saved-log audit: exactly one initial reset, no internal or pending resets, monotone time with 0.02~s microsteps (0.04~s environment steps), constant targets, state-chain consistency, and recomputed predicates and OR/END labels. Across \RfourMicroCount{} microsteps every check has zero mismatches. The \RfourPrefixCount{} actions before first success (\RfourPrefixHitCount{} including its action) are bitwise identical to the sealed original actions. Old/new any-step outcomes have zero mismatches per case, stream and checkpoint. These full-grid checks do not rerun physics. Patch diffs, the environment freeze and a minimal saved-log verification entry point are archived at Zenodo, DOI \href{https://doi.org/10.5281/zenodo.23350373}{10.5281/zenodo.23350373}, and mirrored on GitHub.

\input{tables/table_reacher_termination_contrasts.tex}
\input{tables/table_reacher_microstep_end.tex}
\input{tables/table_reacher_correction_timeline.tex}


\FloatBarrier
\section{Compute accounting}
The formal runs used the recorded GPU-hour totals in the execution receipts. Compute accounting is reported for reproducibility and planning; it does not enter any success contrast. The analysis, figure, and table generators run on CPU from the machine-readable result files.
\input{tables/table_compute_combined.tex}

\FloatBarrier
\printbibliography
\end{document}
